\chapter{Programación Dinámica e Interacción Política Generalizada}
\label{cap:dp}

\section{El Problema del Control Óptimo}
En el capítulo anterior aprendimos a resolver el problema de la \textbf{evaluación de predicción}: dada una política fija $\pi$, calcular su función de valor $V^\pi(s)$. Sin embargo, un agente inteligente autónomo no debe limitarse a resignarse ante una política arbitraria; debe descubrir la mejor política posible.

Decimos que una política $\pi^*$ es \textbf{óptima} si su valor esperado acumulado es mayor o igual que el de cualquier otra política para todos y cada uno de los estados del entorno:
\begin{equation}
  V^{\pi^*}(s) \ge V^\pi(s) \quad \forall s \in \mathcal{S}, \quad \forall \pi
\end{equation}
Todas las políticas óptimas comparten la misma función de valor óptima de estado, denotada $V^*(s) \doteq \max_\pi V^\pi(s)$, y la misma función de valor de acción óptima $Q^*(s,a) \doteq \max_\pi Q^\pi(s,a)$.

\section{Ecuaciones de Bellman de Optimalidad}
Bajo una política óptima, la decisión tomada en cada estado debe coincidir con la acción que maximiza la suma de la recompensa esperada inmediata más el valor óptimo del estado sucesor descontado:
\begin{align}
  V^*(s) &= \max_{a \in \mathcal{A}} Q^*(s, a) \notag \\
  &= \max_{a \in \mathcal{A}} \left[ \mathcal{R}(s, a) + \gamma \sum_{s' \in \mathcal{S}} \mathcal{P}_{ss'}^a V^*(s') \right]
\end{align}
De forma complementaria, para el valor de acción óptimo:
\begin{equation}
  Q^*(s, a) = \mathcal{R}(s, a) + \gamma \sum_{s' \in \mathcal{S}} \mathcal{P}_{ss'}^a \max_{a' \in \mathcal{A}} Q^*(s', a')
\end{equation}

\begin{ideaclave}
Nótese una diferencia matemática crítica: la Ecuación de Expectativa de Bellman era lineal, lo que permitía resolverla mediante inversión matricial. Por el contrario, la Ecuación de Optimalidad contiene el operador no lineal $\max_{a}$, lo que invalida el álgebra lineal directa y nos obliga a recurrir a algoritmos iterativos de punto fijo.
\end{ideaclave}

\section{Policy Iteration y Value Iteration}
La Programación Dinámica (DP) ofrece dos algoritmos fundamentales para hallar $V^*$ y $\pi^*$, asumiendo que conocemos perfectamente el modelo del entorno ($\mathcal{P}$ y $\mathcal{R}$):

\subsection{1. Policy Iteration (Iteración de Políticas)}
Comienza con una política arbitraria $\pi_0$ y repite dos etapas consecutivas:
\begin{enumerate}
  \item \textbf{Evaluación de Política:} Resuelve iterativamente $V^{\pi_k}$ hasta convergencia:
  \[
    V_{m+1}(s) \leftarrow \sum_{a} \pi_k(a \mid s) \left[ \mathcal{R}(s,a) + \gamma \sum_{s'} \mathcal{P}_{ss'}^a V_m(s') \right]
  \]
  \item \textbf{Mejora de Política:} Construye una política estrictamente mejor eligiendo de forma codiciosa (\textit{greedy}) con respecto al valor calculado:
  \[
    \pi_{k+1}(s) \leftarrow \arg\max_{a \in \mathcal{A}} \left[ \mathcal{R}(s,a) + \gamma \sum_{s'} \mathcal{P}_{ss'}^a V^{\pi_k}(s') \right]
  \]
\end{enumerate}
El \textbf{Teorema de Mejora de Política} garantiza que $V^{\pi_{k+1}}(s) \ge V^{\pi_k}(s)$ para todo estado $s$. Al existir un número finito de políticas deterministas en un MDP finito, el algoritmo converge exactamente a $\pi^*$ en un número finito de iteraciones.

\subsection{2. Value Iteration (Iteración de Valores)}
La evaluación completa de la política hasta su estricta convergencia en cada ciclo es computacionalmente costosa. Richard Bellman propuso truncar la evaluación a una sola pasada, incorporando la actualización de optimalidad directamente:
\begin{equation}
  V_{k+1}(s) \leftarrow \max_{a \in \mathcal{A}} \left[ \mathcal{R}(s, a) + \gamma \sum_{s' \in \mathcal{S}} \mathcal{P}_{ss'}^a V_k(s') \right]
\end{equation}
Por el \textbf{Teorema del Punto Fijo de Banach}, el operador de optimalidad de Bellman es una $\gamma$-contracción estricta en la norma del supremo ($\|T(u) - T(v)\|_\infty \le \gamma \|u - v\|_\infty$). Por tanto, partiendo de cualquier vector inicial $V_0$, la sucesión $\{V_k\}$ converge invariablemente a la solución única $V^*$.

\section{Interacción Política Generalizada (GPI)}
Casi todos los métodos de Aprendizaje por Refuerzo pueden conceptualizarse como instancias de \textbf{Interacción Política Generalizada} (\textit{Generalized Policy Iteration}, GPI). En GPI interactúan dos procesos simultáneos:
\begin{itemize}
  \item El proceso de evaluación busca que la función de valor sea consistente con la política actual.
  \item El proceso de mejora busca que la política sea codiciosa respecto a la función de valor actual.
\end{itemize}
Ambos procesos cooperan y compiten dinámicamente hasta que alcanzan un estado de equilibrio en el que la política es codiciosa con respecto a su propia función de valor, lo que sólo ocurre en $(\pi^*, V^*)$.

\begin{optibot}
El fabricante entrega a los ingenieros los planos completos del almacén: una cuadrícula de $3 \times 3$ celdas donde se conocen con absoluta exactitud las probabilidades de movimiento de OptiBot y las penalizaciones por paso ($-1$ en cada movimiento, $+10$ al llegar al punto de entrega en la celda $(2,2)$). Con esta información perfecta, OptiBot no necesita experimentar a ciegas: puede calcular la política de navegación óptima directamente en su procesador central mediante Value Iteration.
\end{optibot}

\section{Implementación en Python: Value Iteration para Cuadrícula}
\begin{lstlisting}[style=rlcode]
import numpy as np

def value_iteration(grid_size=3, gamma=0.9, theta=1e-4):
    V = np.zeros((grid_size, grid_size))
    target = (grid_size - 1, grid_size - 1)
    actions = [(-1, 0), (1, 0), (0, -1), (0, 1)] # Arriba, Abajo, Izq, Der

    while True:
        delta = 0
        for r in range(grid_size):
            for c in range(grid_size):
                if (r, c) == target:
                    continue
                v_old = V[r, c]
                q_candidates = []
                for dr, dc in actions:
                    nr, nc = max(0, min(grid_size-1, r+dr)), max(0, min(grid_size-1, c+dc))
                    reward = 10.0 if (nr, nc) == target else -1.0
                    q_candidates.append(reward + gamma * V[nr, nc])
                V[r, c] = max(q_candidates)
                delta = max(delta, abs(v_old - V[r, c]))
        if delta < theta:
            break
    return V

V_optimo = value_iteration()
print("Mapa de Valores Optimos V*(s):")
print(np.round(V_optimo, 2))
\end{lstlisting}

\begin{notabox}
\textbf{La limitación que motiva el siguiente capítulo:} La Programación Dinámica es conceptualmente perfecta pero padece una grave limitación práctica: \textbf{requiere un modelo completo y perfecto de la física del entorno} ($\mathcal{P}$ y $\mathcal{R}$). En el mundo real, los almacenes cambian, los suelos se desgastan y nadie conoce las probabilidades exactas de transición. Necesitamos algoritmos capaces de aprender \textbf{sin modelo}, basándose exclusivamente en la experiencia empírica. Ese es el reino de los \textbf{Métodos Monte Carlo}.
\end{notabox}
